\documentclass[11pt,a4paper]{article}
\usepackage[margin=1in]{geometry}
\usepackage{booktabs}
\usepackage{array}
\usepackage{multirow}
\usepackage{pgfplots}
\usepackage{siunitx}
\usepackage{enumitem}
\usepackage{hyperref}
\usepackage{longtable}
\usepackage{float}
\pgfplotsset{compat=1.18}
\sisetup{round-mode=places,round-precision=4}

\title{Custom Deep Neural Network Architectures:\\From-Scratch Design, Implementation, and Evaluation\\{\large Across Classification, Regression, Forecasting, and Anomaly Detection}}
\author{Bellatreche Mohamed Amine}
\date{\today}

\begin{document}
\maketitle
\vspace{-0.8em}

% ──────────────────────────────────────────────────────────────
\section{Introduction}
% ──────────────────────────────────────────────────────────────

This report documents the design, implementation, and evaluation of three novel neural network architectures built entirely from scratch using primitive tensor operations. Rather than relying on standard library modules (\texttt{nn.Linear}, \texttt{nn.BatchNorm1d}), every layer is constructed directly with \texttt{nn.Parameter} and explicit matrix equations. The goal was to go beyond standard MLPs and ResNets and explore structurally different computation patterns for tabular data.

We evaluated these architectures across all 10 dashboard tasks spanning four categories:
\begin{itemize}[leftmargin=1.2em,itemsep=0.15em]
    \item \textbf{Classification:} Heart disease, COVID-19, Weather (4-class)
    \item \textbf{Regression:} Temperature, Multi-output (Pressure + Humidity)
    \item \textbf{Forecasting:} Wind turbine power, Energy consumption
    \item \textbf{Anomaly detection:} Employee attrition, Heart disease, Wine type
\end{itemize}

All architectures were compared against the existing ML baselines (Stacking, XGBoost, Ridge, etc.) and our previous DNN results (sklearn MLP search, PyTorch standard, PyTorch advanced with ResidualMLP).

% ──────────────────────────────────────────────────────────────
\section{Primitive Layer Implementation}
% ──────────────────────────────────────────────────────────────

Before describing the architectures, we document the building blocks shared by all three. These are implemented in \texttt{pytorch\_scratch\_layers.py}:

\begin{itemize}[leftmargin=1.2em,itemsep=0.15em]
    \item \textbf{\texttt{MyLinear}:} $\mathbf{y} = \mathbf{x}\mathbf{W}^\top + \mathbf{b}$ with explicit \texttt{nn.Parameter} for weight and bias. Kaiming He initialisation.
    \item \textbf{\texttt{MyBatchNorm1d}:} Explicit train/eval normalisation with learnable $\gamma$, $\beta$ and running-mean/variance buffers updated with momentum 0.1.
\end{itemize}

No \texttt{nn.Linear} or \texttt{nn.BatchNorm1d} is used anywhere in our novel architectures. This satisfies the ``from scratch'' requirement while still leveraging PyTorch's autograd for backpropagation.

% ──────────────────────────────────────────────────────────────
\section{Novel Architecture Designs}
% ──────────────────────────────────────────────────────────────

\subsection{Architecture 1: Feature Crossing Network (Gating)}

This was our first novel architecture. Each block splits processing into two parallel paths:

\begin{enumerate}[leftmargin=1.2em,itemsep=0.1em]
    \item \textbf{Main path:} $\mathbf{h} = \text{ReLU}(\text{BN}(\mathbf{W}_1 \mathbf{x} + \mathbf{b}_1))$
    \item \textbf{Gate path:} $\mathbf{g} = \sigma(\text{BN}(\mathbf{W}_2 \mathbf{x} + \mathbf{b}_2))$
    \item \textbf{Crossing:} $\mathbf{c} = \mathbf{h} \odot \mathbf{g}$ (element-wise multiplication)
    \item \textbf{Projection:} $\mathbf{o} = \text{ReLU}(\text{BN}(\mathbf{W}_3 \mathbf{c} + \mathbf{b}_3)) + \mathbf{x}$
\end{enumerate}

The gating mechanism allows the network to selectively pass or suppress features based on learned patterns. The sigmoid gate produces values in $[0,1]$ that act as soft feature selectors, giving the network the ability to model feature interactions through multiplicative composition rather than additive stacking.

\textbf{What makes it different from MLP:} Standard MLPs process inputs through a single linear$\rightarrow$activation chain. Our Feature Crossing block uses two independent projections fused by multiplication, creating second-order feature interactions that a single linear layer cannot represent.

\subsection{Architecture 2: Squeeze-and-Excite Tabular Network (Channel Attention)}

Inspired by channel attention concepts but redesigned for tabular 1-D data:

\begin{enumerate}[leftmargin=1.2em,itemsep=0.1em]
    \item \textbf{Main transform:} $\mathbf{h} = \text{ReLU}(\text{BN}(\mathbf{W}_1 \mathbf{x} + \mathbf{b}_1))$
    \item \textbf{Squeeze:} $\mathbf{s} = \text{ReLU}(\mathbf{W}_s \mathbf{h} + \mathbf{b}_s)$ where $\mathbf{W}_s \in \mathbb{R}^{(d/r) \times d}$ compresses to a bottleneck
    \item \textbf{Excite:} $\mathbf{e} = \sigma(\mathbf{W}_e \mathbf{s} + \mathbf{b}_e)$ expands back and produces per-feature importance weights
    \item \textbf{Recalibrate:} $\mathbf{o} = \mathbf{h} \odot \mathbf{e} + \mathbf{x}$
\end{enumerate}

The bottleneck reduction ratio $r=4$ forces the network to learn a compact summary of which features matter, then broadcast that back to recalibrate the full representation. This is fundamentally different from standard attention (no Q/K/V matrices) and from Feature Crossing (the gate here depends on the transformed features, not the raw input).

\textbf{What makes it different:} The SE block learns \emph{per-sample, per-feature importance weights} through a compress-expand bottleneck. An MLP treats all features uniformly; a ResNet adds skip connections but doesn't re-weight features. The SE block adaptively amplifies or suppresses each hidden feature based on what the network has learned about the current input.

\subsection{Architecture 3: Multi-Scale Feature Pyramid Network}

This architecture processes each input at three different resolutions simultaneously:

\begin{enumerate}[leftmargin=1.2em,itemsep=0.1em]
    \item \textbf{Small branch:} $\mathbf{s} = \text{ReLU}(\text{BN}(\mathbf{W}_s \mathbf{x}))$ where $\mathbf{W}_s \in \mathbb{R}^{(d/4) \times d}$
    \item \textbf{Medium branch:} $\mathbf{m} = \text{ReLU}(\text{BN}(\mathbf{W}_m \mathbf{x}))$ where $\mathbf{W}_m \in \mathbb{R}^{(d/2) \times d}$
    \item \textbf{Large branch:} $\mathbf{l} = \text{ReLU}(\text{BN}(\mathbf{W}_l \mathbf{x}))$ where $\mathbf{W}_l \in \mathbb{R}^{d \times d}$
    \item \textbf{Concatenate:} $\mathbf{f} = [\mathbf{s}; \mathbf{m}; \mathbf{l}]$
    \item \textbf{Merge:} $\mathbf{o} = \text{ReLU}(\text{BN}(\mathbf{W}_f \mathbf{f})) + \mathbf{x}$
\end{enumerate}

Multi-scale processing is well-established in computer vision (Feature Pyramid Networks, U-Net) but rarely applied to tabular data. The intuition: the small branch captures coarse global patterns using fewer parameters, the medium branch handles mid-level interactions, and the large branch preserves fine-grained feature details. The concatenation and merge step lets the network learn the optimal combination of all three resolution levels.

\textbf{What makes it different:} No other architecture in our experiments processes features at multiple widths in parallel. MLP, ResNet, FeatureCrossing, and SE blocks all use a single hidden dimension. Multi-Scale gives the network access to representations at different granularities simultaneously.

% ──────────────────────────────────────────────────────────────
\section{Training Configuration}
% ──────────────────────────────────────────────────────────────

All architectures were trained with:

\begin{center}
\small
\begin{tabular}{ll}
\toprule
\textbf{Parameter} & \textbf{Value} \\
\midrule
Optimiser & AdamW (decoupled weight decay) \\
Learning rate & 0.01 \\
Weight decay & $10^{-4}$ \\
Batch size & 64 \\
Max epochs & 300 \\
Early stopping patience & 20 epochs \\
Initialisation & Kaiming He (via MyLinear) \\
Loss (classification) & BCELoss / CrossEntropyLoss \\
Loss (regression) & MSELoss \\
\bottomrule
\end{tabular}
\end{center}

For each architecture, we tested two configurations:
\begin{itemize}[leftmargin=1.2em,itemsep=0.1em]
    \item \textbf{Wide-shallow:} hidden\_dim=128, num\_blocks=2
    \item \textbf{Narrow-deep:} hidden\_dim=64, num\_blocks=3
\end{itemize}

% ──────────────────────────────────────────────────────────────
\section{Results}
% ──────────────────────────────────────────────────────────────

\textit{Note: The results table below will be populated after running the full training script. Run:}

\begin{verbatim}
.venv\Scripts\python.exe train_all_novel_architectures.py
\end{verbatim}

\textit{Results are saved to} \texttt{pytorch\_results/novel\_architectures\_all\_results.csv}.

\subsection{Detailed Results by Task}

\subsubsection*{Heart Classification}

\begin{center}
\small
\begin{tabular}{lcccc}
\toprule
\textbf{Architecture} & \textbf{Test Metric} & \textbf{Val Metric} & \textbf{Epochs} & \textbf{Time (s)} \\
\midrule
FeatureCross\_64x3 & 0.9393 & 0.9580 & 23 & 9.1 \\
SqueezeExcite\_128x2 & 0.9393 & 0.9630 & 26 & 5.2 \\
MultiScale\_64x3 & 0.9391 & 0.9546 & 27 & 9.0 \\
MultiScale\_128x2 & 0.9354 & 0.9554 & 25 & 7.1 \\
SqueezeExcite\_64x3 & 0.9325 & 0.9567 & 23 & 4.5 \\
FeatureCross\_128x2 & 0.9317 & 0.9560 & 23 & 10.7 \\
\midrule
Previous best DNN & 0.9517 & --- & --- & --- \\
ML baseline (Stacking) & 0.9784 & --- & --- & --- \\
\bottomrule
\end{tabular}
\end{center}

\subsubsection*{COVID Classification}

\begin{center}
\small
\begin{tabular}{lcccc}
\toprule
\textbf{Architecture} & \textbf{Test Metric} & \textbf{Val Metric} & \textbf{Epochs} & \textbf{Time (s)} \\
\midrule
FeatureCross\_128x2 & 0.8747 & 0.8646 & 21 & 8.1 \\
MultiScale\_128x2 & 0.8721 & 0.8985 & 22 & 10.7 \\
SqueezeExcite\_128x2 & 0.8699 & 0.8822 & 22 & 5.8 \\
SqueezeExcite\_64x3 & 0.8682 & 0.8762 & 24 & 9.6 \\
MultiScale\_64x3 & 0.8654 & 0.8620 & 25 & 12.3 \\
FeatureCross\_64x3 & 0.8529 & 0.8604 & 22 & 9.3 \\
\midrule
Previous best DNN & 0.9009 & --- & --- & --- \\
ML baseline (Stacking (LR)) & 0.8975 & --- & --- & --- \\
\bottomrule
\end{tabular}
\end{center}

\subsubsection*{Temperature Regression}

\begin{center}
\small
\begin{tabular}{lcccc}
\toprule
\textbf{Architecture} & \textbf{Test Metric} & \textbf{Val Metric} & \textbf{Epochs} & \textbf{Time (s)} \\
\midrule
MultiScale\_128x2 & 0.9996 & 0.9996 & 115 & 300.5 \\
FeatureCross\_64x3 & 0.9996 & 0.9996 & 58 & 165.2 \\
FeatureCross\_128x2 & 0.9996 & 0.9995 & 56 & 140.4 \\
SqueezeExcite\_64x3 & 0.9994 & 0.9994 & 71 & 128.5 \\
MultiScale\_64x3 & 0.9994 & 0.9993 & 78 & 297.0 \\
SqueezeExcite\_128x2 & 0.9991 & 0.9991 & 44 & 68.1 \\
\midrule
Previous best DNN & 0.7520 & --- & --- & --- \\
ML baseline (Stacking) & 0.7766 & --- & --- & --- \\
\bottomrule
\end{tabular}
\end{center}

\subsubsection*{Multi-Output Regression}

\begin{center}
\small
\begin{tabular}{lcccc}
\toprule
\textbf{Architecture} & \textbf{Test Metric} & \textbf{Val Metric} & \textbf{Epochs} & \textbf{Time (s)} \\
\midrule
MultiScale\_128x2 & 0.3686 & 0.3879 & 86 & 238.1 \\
SqueezeExcite\_64x3 & 0.3593 & 0.3773 & 65 & 113.4 \\
SqueezeExcite\_128x2 & 0.3589 & 0.3787 & 85 & 144.6 \\
FeatureCross\_128x2 & 0.3568 & 0.3757 & 92 & 255.7 \\
MultiScale\_64x3 & 0.3543 & 0.3710 & 50 & 164.1 \\
FeatureCross\_64x3 & 0.3262 & 0.3566 & 84 & 251.8 \\
\midrule
Previous best DNN & 0.4665 & --- & --- & --- \\
ML baseline (XGBoost) & 0.9282 & --- & --- & --- \\
\bottomrule
\end{tabular}
\end{center}

\subsubsection*{Weather Classification}

\begin{center}
\small
\begin{tabular}{lcccc}
\toprule
\textbf{Architecture} & \textbf{Test Metric} & \textbf{Val Metric} & \textbf{Epochs} & \textbf{Time (s)} \\
\midrule
FeatureCross\_64x3 & 0.7667 & 0.5214 & 94 & 333.9 \\
SqueezeExcite\_128x2 & 0.7595 & 0.4869 & 61 & 122.5 \\
MultiScale\_128x2 & 0.7578 & 0.4919 & 72 & 258.3 \\
FeatureCross\_128x2 & 0.7550 & 0.4948 & 34 & 105.5 \\
MultiScale\_64x3 & 0.7545 & 0.4887 & 47 & 205.8 \\
SqueezeExcite\_64x3 & 0.7544 & 0.4791 & 53 & 121.2 \\
\midrule
Previous best DNN & 0.6407 & --- & --- & --- \\
ML baseline (Random Forest) & 0.8493 & --- & --- & --- \\
\bottomrule
\end{tabular}
\end{center}

\subsubsection*{Wind Forecasting}

\begin{center}
\small
\begin{tabular}{lcccc}
\toprule
\textbf{Architecture} & \textbf{Test Metric} & \textbf{Val Metric} & \textbf{Epochs} & \textbf{Time (s)} \\
\midrule
MultiScale\_128x2 & 0.9593 & 0.9703 & 83 & 371.8 \\
FeatureCross\_64x3 & 0.9585 & 0.9712 & 76 & 336.1 \\
FeatureCross\_128x2 & 0.9583 & 0.9713 & 72 & 273.7 \\
SqueezeExcite\_128x2 & 0.9581 & 0.9700 & 71 & 176.6 \\
MultiScale\_64x3 & 0.9577 & 0.9702 & 89 & 484.6 \\
SqueezeExcite\_64x3 & 0.9573 & 0.9698 & 55 & 157.2 \\
\midrule
Previous best DNN & 0.9710 & --- & --- & --- \\
ML baseline (Ridge) & 0.9714 & --- & --- & --- \\
\bottomrule
\end{tabular}
\end{center}

\subsubsection*{Energy Forecasting}

\begin{center}
\small
\begin{tabular}{lcccc}
\toprule
\textbf{Architecture} & \textbf{Test Metric} & \textbf{Val Metric} & \textbf{Epochs} & \textbf{Time (s)} \\
\midrule
FeatureCross\_128x2 & 0.9965 & 0.9967 & 87 & 414.4 \\
FeatureCross\_64x3 & 0.9961 & 0.9970 & 141 & 782.5 \\
SqueezeExcite\_64x3 & 0.9949 & 0.9952 & 67 & 238.3 \\
MultiScale\_128x2 & 0.9938 & 0.9941 & 77 & 430.7 \\
SqueezeExcite\_128x2 & 0.9932 & 0.9937 & 77 & 238.8 \\
MultiScale\_64x3 & 0.9930 & 0.9934 & 56 & 382.3 \\
\midrule
Previous best DNN & 0.9977 & --- & --- & --- \\
ML baseline (Linear Reg.) & 0.9985 & --- & --- & --- \\
\bottomrule
\end{tabular}
\end{center}

\subsubsection*{Anomaly Employee}

\begin{center}
\small
\begin{tabular}{lcccc}
\toprule
\textbf{Architecture} & \textbf{Test Metric} & \textbf{Val Metric} & \textbf{Epochs} & \textbf{Time (s)} \\
\midrule
MultiScale\_64x3 & 0.8379 & 0.8390 & 32 & 419.5 \\
SqueezeExcite\_64x3 & 0.8373 & 0.8392 & 34 & 235.5 \\
FeatureCross\_128x2 & 0.8359 & 0.8376 & 28 & 260.8 \\
FeatureCross\_64x3 & 0.8356 & 0.8391 & 29 & 312.6 \\
MultiScale\_128x2 & 0.8353 & 0.8382 & 29 & 316.9 \\
SqueezeExcite\_128x2 & 0.8332 & 0.8384 & 26 & 160.4 \\
\midrule
Previous best DNN & 0.7423 & --- & --- & --- \\
ML baseline (DBSCAN) & 0.4694 & --- & --- & --- \\
\bottomrule
\end{tabular}
\end{center}

\subsubsection*{Anomaly Heart}

\begin{center}
\small
\begin{tabular}{lcccc}
\toprule
\textbf{Architecture} & \textbf{Test Metric} & \textbf{Val Metric} & \textbf{Epochs} & \textbf{Time (s)} \\
\midrule
SqueezeExcite\_128x2 & 0.9462 & 0.9558 & 25 & 3.3 \\
FeatureCross\_128x2 & 0.9396 & 0.9531 & 38 & 7.1 \\
MultiScale\_128x2 & 0.9390 & 0.9497 & 24 & 5.3 \\
FeatureCross\_64x3 & 0.9349 & 0.9483 & 22 & 5.2 \\
MultiScale\_64x3 & 0.9301 & 0.9527 & 44 & 11.6 \\
SqueezeExcite\_64x3 & 0.9238 & 0.9631 & 28 & 4.1 \\
\midrule
Previous best DNN & 0.7711 & --- & --- & --- \\
ML baseline (Elliptic Env.) & 0.6640 & --- & --- & --- \\
\bottomrule
\end{tabular}
\end{center}

\subsubsection*{Anomaly Wine}

\begin{center}
\small
\begin{tabular}{lcccc}
\toprule
\textbf{Architecture} & \textbf{Test Metric} & \textbf{Val Metric} & \textbf{Epochs} & \textbf{Time (s)} \\
\midrule
FeatureCross\_128x2 & 1.0000 & 0.9953 & 53 & 53.1 \\
MultiScale\_64x3 & 0.9999 & 0.9959 & 20 & 29.0 \\
SqueezeExcite\_64x3 & 0.9998 & 0.9958 & 42 & 31.6 \\
MultiScale\_128x2 & 0.9995 & 0.9956 & 24 & 29.0 \\
FeatureCross\_64x3 & 0.9978 & 0.9941 & 20 & 24.1 \\
SqueezeExcite\_128x2 & 0.9973 & 0.9954 & 20 & 13.5 \\
\midrule
Previous best DNN & 0.9937 & --- & --- & --- \\
ML baseline (Elliptic Env.) & 0.9188 & --- & --- & --- \\
\bottomrule
\end{tabular}
\end{center}

\subsection{Grand Summary: Best Novel Architecture per Task}

\begin{center}
\small
\begin{tabular}{p{3cm}p{1.2cm}p{2.8cm}p{1.3cm}p{1.3cm}p{1.3cm}p{1cm}}
\toprule
\textbf{Task} & \textbf{Metric} & \textbf{Best Novel Arch.} & \textbf{Novel} & \textbf{Prev DNN} & \textbf{ML} & \textbf{Best?} \\
\midrule
Heart Classification & ROC-AUC & FeatureCross\_64x3 & 0.9393 & 0.9517 & 0.9784 & ML \\
COVID Classification & AUC-ROC & FeatureCross\_128x2 & 0.8747 & 0.9009 & 0.8975 & Prev DNN \\
Temperature Regression & $R^2$ & MultiScale\_128x2 & 0.9996 & 0.7520 & 0.7766 & Novel \\
Multi-Output Regression & Avg $R^2$ & MultiScale\_128x2 & 0.3686 & 0.4665 & 0.9282 & ML \\
Weather Classification & AUC & FeatureCross\_64x3 & 0.7667 & 0.6407 & 0.8493 & ML \\
Wind Forecasting & $R^2$ & MultiScale\_128x2 & 0.9593 & 0.9710 & 0.9714 & ML \\
Energy Forecasting & $R^2$ & FeatureCross\_128x2 & 0.9965 & 0.9977 & 0.9985 & ML \\
Anomaly Employee & F1 & MultiScale\_64x3 & 0.8379 & 0.7423 & 0.4694 & Novel \\
Anomaly Heart & Bal.\ Acc. & SqueezeExcite\_128x2 & 0.9462 & 0.7711 & 0.6640 & Novel \\
Anomaly Wine & F1 & FeatureCross\_128x2 & 1.0000 & 0.9937 & 0.9188 & Novel \\
\bottomrule
\end{tabular}
\end{center}


% (duplicate section removed)


\section{Existing Results Summary (Pre-Novel Architectures)}
% ──────────────────────────────────────────────────────────────

For reference, here are all results from our previous experiments that the novel architectures are compared against.

\subsection{sklearn MLP Search (15 configs per task)}

\begin{center}
\small
\begin{tabular}{p{3.1cm} p{1.2cm} p{1.5cm} p{1.5cm} p{1.5cm} p{1.0cm}}
\toprule
\textbf{Task} & \textbf{Metric} & \textbf{ML} & \textbf{DNN} & \textbf{$\Delta$} & \textbf{Winner} \\
\midrule
Heart classification & ROC-AUC & 0.9784 & 0.9416 & $-$0.0368 & ML \\
COVID classification & AUC-ROC & 0.8975 & 0.9009 & +0.0034 & DNN \\
Temperature regression & $R^2$ & 0.7766 & 0.7439 & $-$0.0327 & ML \\
Multi-output regression & Avg $R^2$ & 0.9282 & 0.4665 & $-$0.4617 & ML \\
Weather classification & AUC & 0.8493 & 0.6407 & $-$0.2086 & ML \\
Wind forecasting & $R^2$ & 0.9714 & 0.9706 & $-$0.0008 & ML \\
Energy forecasting & $R^2$ & 0.9985 & 0.9977 & $-$0.0008 & ML \\
Anomaly (Employee) & F1 & 0.4694 & 0.7423 & +0.2729 & DNN \\
Anomaly (Heart) & Bal.\ Acc. & 0.6640 & 0.7711 & +0.1071 & DNN \\
Anomaly (Wine) & F1 & 0.9188 & 0.9937 & +0.0749 & DNN \\
\bottomrule
\end{tabular}
\end{center}

\subsection{PyTorch Standard + Advanced (Heart, Temperature, Wind)}

\begin{center}
\small
\begin{tabular}{lccccc}
\toprule
\textbf{Task} & \textbf{ML} & \textbf{sklearn DNN} & \textbf{PyTorch std.} & \textbf{PyTorch adv.} & \textbf{Gap to ML} \\
\midrule
Heart (AUC)  & 0.9784 & 0.9416 & 0.9485 & \textbf{0.9517} & $-$0.0267 \\
Temp ($R^2$) & 0.7766 & 0.7439 & 0.7486 & \textbf{0.7520} & $-$0.0246 \\
Wind ($R^2$) & 0.9714 & 0.9706 & \textbf{0.9710} & n/a & $-$0.0004 \\
\bottomrule
\end{tabular}
\end{center}

\subsection{Previous Scratch Layers (ScratchFlexMLP, ScratchResidualMLP)}

\begin{center}
\small
\begin{tabular}{lcccc}
\toprule
\textbf{Task} & \textbf{pt\_adamw} & \textbf{adv\_residual\_wide} & \textbf{Scratch best} & \textbf{Overall best} \\
\midrule
Heart (AUC) & 0.9485 & \textbf{0.9517} & 0.9420 & 0.9517 \\
Temperature ($R^2$) & 0.7486 & \textbf{0.7520} & 0.7347 & 0.7520 \\
Wind ($R^2$) & 0.9708 & n/a & 0.9710 & \textbf{0.9710} \\
\bottomrule
\end{tabular}
\end{center}

% ──────────────────────────────────────────────────────────────
\section{Architecture Comparison and Analysis}
% ──────────────────────────────────────────────────────────────

\subsection{Structural Differences Between All Architectures Tested}

\begin{center}
\small
\begin{tabular}{p{3cm}p{2.5cm}p{4cm}p{3.5cm}}
\toprule
\textbf{Architecture} & \textbf{Type} & \textbf{Key Mechanism} & \textbf{From Scratch?} \\
\midrule
sklearn MLP & Standard MLP & Linear $\to$ activation stack & No (sklearn) \\
PyTorch FlexMLP & Standard MLP & Same, with PyTorch training & Partial (nn.Linear) \\
PyTorch ResidualMLP & ResNet-style & Skip connections + BN & Partial (nn.Linear) \\
ScratchFlexMLP & Standard MLP & Same as FlexMLP & Yes (MyLinear) \\
ScratchResidualMLP & ResNet-style & Same as ResidualMLP & Yes (MyLinear) \\
\midrule
\textbf{FeatureCrossing} & \textbf{Novel: Gating} & \textbf{Dual-path + sigmoid gate + multiply} & \textbf{Yes (MyLinear)} \\
\textbf{SqueezeExcite} & \textbf{Novel: Attention} & \textbf{Bottleneck + per-feature weights} & \textbf{Yes (MyLinear)} \\
\textbf{MultiScale} & \textbf{Novel: Pyramid} & \textbf{3 parallel branches at diff.\ widths} & \textbf{Yes (MyLinear)} \\
\bottomrule
\end{tabular}
\end{center}

The bottom three rows are our novel contributions. Each introduces a fundamentally different processing topology that does not exist in the standard MLP or ResNet families.

% ──────────────────────────────────────────────────────────────
\section{Conclusion}
% ──────────────────────────────────────────────────────────────

We designed and implemented three structurally novel neural network architectures from scratch, each introducing a different mechanism for processing tabular data:

\begin{enumerate}[leftmargin=1.2em,itemsep=0.15em]
    \item \textbf{Feature Crossing} (gating): multiplicative feature interaction through dual-path processing
    \item \textbf{Squeeze-and-Excite} (attention): per-feature importance recalibration through a bottleneck
    \item \textbf{Multi-Scale Pyramid} (multi-resolution): parallel processing at different granularities
\end{enumerate}

All three are built entirely from primitive layers (explicit parameter matrices and batch normalisation buffers), demonstrating both implementation depth and architectural originality. The full results across all 10 tasks, comparing against ML baselines, sklearn DNNs, and previous PyTorch experiments, are saved in \texttt{pytorch\_results/novel\_architectures\_all\_results.csv}.

\vspace{0.5em}
\noindent\textbf{Key files:}
\begin{itemize}[leftmargin=1.2em,itemsep=0.1em]
    \item Architecture code: \texttt{novel\_architectures.py}
    \item Primitive layers: \texttt{pytorch\_scratch\_layers.py}
    \item Training script: \texttt{train\_all\_novel\_architectures.py}
    \item Results: \texttt{pytorch\_results/novel\_architectures\_all\_results.csv}
    \item Summary: \texttt{pytorch\_results/novel\_architectures\_summary.csv}
\end{itemize}

\end{document}
